\section{Step 118 Record lateness once per interval}
Fill \texttt{vp\_late}. This path is reached only when observed time is at or beyond the scheduled deadline. Subtract the deadline to obtain nonnegative lateness. Increment the late-interval count only for a positive difference, and update maximum lateness if that difference is larger.

Count intervals rather than samples. An interval that needed several short waits is still one interval. If an oversleep reaches 1047 for deadline 1040, its lateness is seven milliseconds; the next scheduled deadline remains 1080. Repeated late samples do not justify dropping a frame in this player.

Use \texttt{s118-1} and stop before \texttt{vp\_interval\_complete}.

\paragraph{Observe} An exact 1040 sample records zero lateness. A 1047 sample records one late interval and maximum lateness seven.

\Needspace{8\baselineskip}
\paragraph{Check} How often can one interval increment the late counter?
\begin{enumerate}[label=\Alph*.]
\item Once for every clock sample.
\item Once for every changed pixel.
\item At most once, when that interval completes late.
\end{enumerate}
